\chapter{Critical Review and Reflection}

\section{Achievement of Objectives}

The primary objective of this project was to engineer and evaluate a suite of deep learning models capable of playing chess without search algorithms. The framework implements all six proposed architectures (ConvNet, ResNet, SquareTransformer, PieceTransformer, GCN, GAT), each trainable under identical data conditions with interchangeable output heads. Both the modular backbone-head design and the streaming data pipeline.

The secondary objective of cross-architecture comparison are supported by the benchmarking pipeline, which records and exports PGN files for qualitative analysis.

Beyond the initial architecture, the framework evolved during development. The ResNet backbone was augmented with Squeeze-and-Excitation blocks for dynamic channel attention. A spatial policy head replaced global average pooling to preserve square-level information. The data pipeline was also designed to be capable of handling over one billion positions.
% [figure] Timeline-style diagram showing key engineering milestones across the project.

\section{What Worked Well}

The abstract \texttt{ChessModel} interface made adding new architectures straightforward. Adding the GAT architecture, for instance, meant implementing a single class with two methods; the training loop, loss functions, and evaluation pipeline needed zero modifications. That backbone-head decoupling justified the upfront design effort many times over.

Row-group-level streaming from Parquet files enabled training on datasets larger than available RAM without changing the training loop. This mattered most when scaling from initial prototyping (100K positions) to full-scale runs (tens of millions of positions). Centralising all hyperparameters in a single YAML file eliminated hard-coded constants and made experiment management straightforward. Combined with checkpoint serialisation, any past experiment can be exactly reproduced.

The decision to inherit root evaluations (with sign flips) for PV-unrolled positions rather than re-evaluating each with Stockfish was both practical and accurate. The drift analysis, showing 96\% of inherited values within $\pm 100$ cp, validated this and enabled a large increase in training data diversity at negligible computational cost.

\section{What Didn't Work Well}

The GNN encoder is much slower than the CNN encoder because it must compute attack and defence relations for every square on every board. This makes the GNN data pipeline the bottleneck during training, even though the model itself is no larger than the CNN.
% [figure] Profiling breakdown of data pipeline time spent per encoder type.

The original policy head used global average pooling, discarding spatial information critical for move prediction. This was caught early and replaced with a spatial policy head, but those initial training runs on the pooled representation produced noticeably weaker policy accuracy — confirming that square-level information is essential for move selection.

% TODO: Write more - there were likely other issues worth discussing: debugging GNN edge batching, mixed-precision silent failures on MPS, dataset imbalance problems, or hyperparameter tuning dead ends. Add 2-3 more paragraphs on things that cost you time.

\section{Lessons Learned}

Filtering out low-depth evaluations and balancing the value distribution had a larger impact on model performance than simply increasing dataset size. Data quality consistently mattered more than data quantity.

The choice of board representation (spatial tensor vs \ graph vs \ token sequence) determines what information is accessible to the model and what must be learned from scratch. Encoder design proved as consequential as model architecture.

The ConvNet, despite being the simplest model, was an invaluable debugging tool. Once it was training correctly, issues in more complex architectures (Transformer attention masks, GNN edge batching) could be diagnosed by comparison. Starting with the simplest architecture saved a lot of time.

Mixed-precision training requires platform-specific handling. The asymmetry between CUDA and MPS support for \texttt{GradScaler} required non-obvious conditional logic that initially caused silent numerical instability.

\section{Personal Reflection}

This project deepened my understanding of deep learning engineering well beyond what I covered in lectures. Implementing graph neural networks from scratch — including the non-trivial batching of variable-edge graphs — required engaging with PyTorch Geometric at a level far beyond its tutorial examples. The Transformer implementations reinforced how much positional encoding design matters: absolute, relative, and learnable encodings produced noticeably different training dynamics on the same data.

Managing a multi-architecture codebase also developed my software engineering discipline. The abstract base class pattern — which initially felt like over-engineering for a single-developer project — turned out to be essential when I added the sixth architecture three months into development.

% TODO: Write more - add specific experiences: a moment where something broke badly and how you fixed it, what you would do differently if starting over, how your thinking about neural network design changed over the course of the project. This section should be at least 3-4 paragraphs.

\section{Future Work}

\subsection{Reinforcement Learning}

The framework currently relies entirely on supervised learning from engine-annotated data. Two reinforcement learning paradigms are natural extensions. Recording MCTS visit-count distributions as policy targets and game outcomes as value targets — enabling the model to improve beyond the ceiling of its training data. Alternatively, playing the model against Stockfish at progressively increasing difficulty (varying skill level and search depth) would give a structured reinforcement signal.

% TODO: Write more - which of these two approaches would you try first and why? What compute resources would be needed? What specific challenges do you anticipate?

\subsection{Web Platform}

The planned Python/Svelte web platform would give users a browser-based interface for interacting with trained models. They could play any trained model via a drag-and-drop board interface, or watch two models play each other with spectator support, configurable move delay, and live evaluation bars. A round-robin tournament system — where all trained models play each other with automated Elo ratings calculated from the results — would allow systematic strength comparisons. An analysis board mode for setting up arbitrary positions and comparing move recommendations from different models side by side is also planned.

\subsection{Long-term Research Directions}

Framing chess as a sequence prediction task — where the model autoregressively predicts move continuations from a history of board states, analogous to language modelling — is one direction I find interesting. This Decision Transformer-style approach could enable genuine multi-move planning without explicit search. The backbone--head architecture is also game-agnostic, so adapting the framework to other perfect-information games (e.g., Shogi, Othello) would test the generality of the design.

% TODO: Write more - are there specific papers or ideas that motivate these directions? Which would you prioritise and why?

\section{Broader Impact}

The modular design lowers the barrier to entry for researchers who want to experiment with novel architectures or training regimes without reimplementing the full chess AI stack. The web platform, once complete, would serve as a demonstration tool for communicating deep learning concepts to non-specialist audiences.
